% Clinical Background section. Paste into your own thesis, or \input it.

\section{Geometry as a risk factor}
\label{sec:geometry}

The shape of the coronary tree is a clinical quantity rather than an incidental descriptor, and
the accepted causal chain explains why: geometry and topology determine wall shear stress, wall
shear stress drives the endothelial response, and the endothelial response produces plaque
\cite{rampidis2022geometry}.
Section~\ref{sec:hemodynamics} sets out the mechanics of the first two links; evidence that the
chain holds in patients is the reason particular geometric quantities are worth measuring.

Three geometric characteristics carry prognostic value beyond stenosis severity and plaque
composition: a more proximal lesion location, location at a bifurcation, and increased
tortuosity.
They were identified in patients who subsequently experienced an acute coronary syndrome, against
plaque characteristics measured on the same scans \cite{han2022plaque}.
Each is a property of the tree rather than of the lesion, which licenses treating the geometry of
the coronary tree as an object of study in its own right.

The left main bifurcation angle is the geometric quantity with the firmest evidence behind it.
Juan et al.\ measured it on CCTA in 313 patients grouped by disease severity and found a gradient:
$75.5^\circ \pm 14.7^\circ$ in 40 patients with normal arteries and a coronary calcium score of
zero (a CT-based count of calcified plaque, so zero means none was found),
$81.2^\circ \pm 20.1^\circ$ in 62 with non-obstructive plaque, and $87.3^\circ \pm 18.8^\circ$ in
211 with a stenosis of at least 50\% \cite{juan2017bifurcation}.
Wang et al.\ supply a mechanism for that gradient: they extracted centerline geometry from 76 left
main arteries, clustered it into four morphological phenotypes, and computed time-averaged wall
shear stress for each by computational fluid dynamics, finding the lowest shear around the branch
points of the LAD and most markedly so in the phenotype characterized by a large bifurcation angle
\cite{wang2024lmphenotypes}.
An anatomical gradient and a hemodynamic mechanism pointing the same way are what make left main
bifurcation angle worth extracting here as a feature in its own right, rather than as one of several
junction measurements.
Effect sizes and angle cut-offs differ between studies, so this thesis takes the direction of the
association from the literature but not a threshold.

Tortuosity is the second such quantity.
In 131 patients investigated for suspected coronary artery disease, severe tortuosity --- defined
as three or more consecutive bends of at least $45^\circ$ --- was the factor most strongly and
independently associated with non-obstructive coronary artery disease: patients with it had
roughly eight times the odds of disease (odds ratio 7.96, 95\% confidence interval 1.79 to 35.40)
\cite{zebicmihic2023tortuosity}.
That interval is very wide on 131 patients, so the size of the effect is poorly constrained even
though its presence is not, and the finding is best read as a reason to measure tortuosity rather
than as a coefficient to reproduce.

Dominance completes the set.
Veltman et al.\ followed 1425 patients after CCTA for a median of 24 months, tracking non-fatal
myocardial infarction and death from any cause, and found left dominance carried roughly three
times the risk of one of those events after accounting for other risk factors (adjusted hazard
ratio 3.20, 95\% confidence interval 1.67 to 6.13) \cite{veltman2012dominance}.
The effect was confined to patients who already had significant disease, among whom three-year
event rates were 35\% in left-dominant and 9.5\% in right-dominant circulations; among patients
without significant disease the two groups did not differ \cite{veltman2012dominance}.
Dominance is therefore a modifier of risk in the presence of disease rather than a risk factor on
its own, which is why this thesis treats it as a stratifying variable for the other features rather
than as a predictor to be reported alone.

One reference value resists this treatment.
Curvature is higher in atherosclerotic than in non-atherosclerotic segments, and higher in
end-systole than in diastole \cite{rampidis2022geometry}; the review reporting both does not give
the cohort sizes behind them, so they set a scale rather than a population baseline.
The cardiac-phase dependence matters methodologically regardless: images here are reconstructed in
mid-diastole (Section~\ref{sec:ccta}), which fixes the phase at which every curvature and tortuosity
feature in this thesis is measured, and makes those values comparable to each other but not directly
to values measured in systole.

Angle, tortuosity and dominance each carry outcome evidence of the kind quoted above.
Descriptors of branching pattern --- how many branches a tree carries and how they are arranged ---
do not, to anything like the same degree, which is part of why describing their distribution across
a large cohort is worth doing at all.
